\documentclass[11pt]{article}
\usepackage[T1]{fontenc}\usepackage[utf8]{inputenc}
\usepackage{lmodern,amsmath,amssymb,amsthm,mathtools,microtype,xcolor}
\usepackage[margin=1in]{geometry}
\usepackage[colorlinks=true,linkcolor=blue!50!black,urlcolor=blue!50!black,citecolor=blue!50!black]{hyperref}
\newtheorem{theorem}{Theorem}[section]\newtheorem{proposition}[theorem]{Proposition}
\newtheorem{lemma}[theorem]{Lemma}\newtheorem{corollary}[theorem]{Corollary}
\theoremstyle{remark}\newtheorem{remark}[theorem]{Remark}
\newcommand{\E}{\mathbb E}\newcommand{\Prb}{\mathbb P}\renewcommand{\Pr}{\mathbb P}\newcommand{\R}{\mathbb R}
\newcommand{\norm}[1]{\left\Vert#1\right\Vert}
\setlength{\emergencystretch}{3em}
\hypersetup{pdftitle={Bayesian neural-network learning with non-Gaussian weight priors},pdfauthor={ChatGPT-6 Astra; Kaarel Hanni; Dmitry Vaintrob}}
\title{Bayesian neural-network learning\\with non-Gaussian weight priors}
\author{ChatGPT-6 Astra \and Kaarel H\"anni \and Dmitry Vaintrob\thanks{Authors are listed in alphabetical order. This research draft has not been formalized in Lean.}}
\date{October 18, 2026}
\begin{document}\maketitle
\begin{abstract}
Can a Bayesian neural network learn a small circuit when its independent
weight priors are not Gaussian? We prove that three width-independent
learning guarantees survive for a natural family of non-Gaussian coordinate
laws: even densities $Z^{-1}e^{-V}$, with variance one and curvature bounded
above and below by positive constants. The guarantees cover the strong-prior
limit, taken separately at each fixed width, and the finite-scale regimes
$\gamma>1$ and $\gamma=1$, in every case at every width at least the
target's bounded-weight network
complexity. The sample bound is polynomial in input dimension and target
complexity at each fixed depth, independently of the learner's width.
The model uses tanh hidden layers, a linear readout, an absolute output
margin imposed at every Boolean input before training, and exact
conditioning on noiseless labels. The prediction guarantee controls the
actual posterior, including mixtures: its correct-label probability exceeds
$2/3$ on almost all test inputs.
The proof uses a translation-entropy bound for target mass and concentration
of strongly log-concave measures. At mean-field scaling, a relative Lindeberg
estimate controls log partition functions for finite covariance tests
with polynomial losses;
an ordinary central limit approximation would not provide the needed bound.
We also give a broader quadratic-tail extension of the strong-prior result.
No efficient posterior computation or polynomial dependence on depth is
asserted.
\end{abstract}

\section{Introduction}
A useful sample bound for a large Bayesian neural network should depend on
the size of a small implementation of the target, rather than on the number
of parameters available to the learner. Gaussian weight priors permit such
bounds in several small-scale and mean-field regimes. This paper establishes
that exact Gaussianity is unnecessary for three of those guarantees.

Our main family consists of independent coordinates with a fixed even
density $q(t)=Z^{-1}e^{-V(t)}$, standardized to variance one, where
\begin{equation}\label{eq:curvature}
 V\in C^2(\R),\qquad 0<m\leq V''(t)\leq M<\infty
 \quad(t\in\R).
\end{equation}
The lower bound supplies concentration. The upper bound makes a translated
copy of a small target inexpensive in relative entropy. These are
restrictions on a one-dimensional weight law, rather than restrictions
selected separately for each network function. A genuinely non-Gaussian
example is the variance-standardized law with density proportional to
$\exp\{-t^2/2-b\log\cosh t\}$, for any fixed $b>0$.

We retain the same layer scales and the same learning rule throughout.
Before training, the prior is conditioned to produce outputs of absolute
value greater than $1/2$ at every Boolean input. Training then conditions
on exact agreement with noiseless labels. At each fixed depth, we prove a
sample bound polynomial in dimension and bounded-weight target complexity,
uniformly over every learner width at least that complexity. There are
three regimes: a common scale tending to zero at fixed width; finite-scale
priors with exponent $\gamma>1$; and the mean-field endpoint $\gamma=1$.
The latter two statements require neither a small-scale limit nor an
infinite-width limit.

The three proofs use a shared statistical principle. A finite collection
of simply describable functions contains all but a sufficiently small
amount of prior mass relative to the target. A union bound eliminates its
inaccurate members after polynomially many examples. In the strong-prior
limit, quadratic tail envelopes select functions of bounded weighted
squared norm. For $\gamma>1$, concentration permits a spectral truncation
with the same effect. At $\gamma=1$, the bulk internal matrices cannot be
removed that way. We instead approximate empirical hidden-layer covariances
by controlled exponential tilts, and then realize those covariances by
small finite networks.

The new difficulty at the mean-field endpoint is that non-Gaussian internal
rows are not determined by their covariance. We resolve it with a relative
Lindeberg comparison for log partition functions. Its error is polynomial
in the tilt budget and independent of the number of input points. This
avoids both an exponential dependence on the tilt budget and a factor equal
to the size of the Boolean cube. All supporting compression and posterior
arguments are included below and in the appendices.

\section{Model and guarantees}\label{sec:model}
Let $d\geq1$, $L\geq2$, and $\mathcal X=\{-1,1\}^d$. Write
$\bar x=(1,x_1,\ldots,x_d)$. A depth-$L$ network has $L-1$ hidden layers
of common width $N$ and computes
\[
 h^0(x)=\bar x,\qquad h^s(x)=\tanh(W^s h^{s-1}(x))\ (1\leq s<L),
 \qquad z_w(x)=W^Lh^{L-1}(x).
\]
There are no separate internal biases. For a Boolean function $g$, let
\[
 U_{N,g}=\{w:g(x)z_w(x)>1/2\ \text{for every }x\in\mathcal X\},
 \qquad U_N=\bigcup_g U_{N,g}.
\]
If $\mu$ is the independent-coordinate parameter prior, its induced
function prior is $\pi(g)=\mu(U_{N,g})/\mu(U_N)$. All priors used below
have positive densities, so any nonempty strict-margin region has positive
mass. In particular the learner is defined whenever the target has a
finite implementation at the chosen width.

Let $C_L(f)$ be the least integer $C\geq1$ for which $f$ has a
width-$C$, depth-$L$ implementation with the stated strict margin and all
weights at most $C$ in absolute value. This combines width and weight
magnitude at the specified depth; it is not an unrestricted gate count.
Given independent inputs $x_1,\ldots,x_n\sim D$ and labels $f(x_i)$,
the posterior is
\[
 \Pi(g\mid S_n)=
 \frac{\pi(g)\mathbf1\{g(x_i)=f(x_i)\ \forall i\}}
 {\sum_h\pi(h)\mathbf1\{h(x_i)=f(x_i)\ \forall i\}}.
\]
Define its Gibbs risk and confidence-failure probability by
\[
 R_D(\Pi)=\E_{g\sim\Pi}\Pr_{x\sim D}\{g(x)\ne f(x)\},\qquad
 F_q(\Pi)=\Pr_{x\sim D}\{\Pi(g(x)=f(x)\mid S_n)\leq q\}.
\]
For every posterior,
\begin{equation}\label{eq:confidence}
 F_q(\Pi)\leq\frac{R_D(\Pi)}{1-q}.
\end{equation}
Thus the statements below control the entire posterior, rather than a
selected minimizer or the mean of its real-valued outputs. Posterior-majority
error is at most $2R_D(\Pi)$, including arbitrary tie breaking.

Let $\xi$ have density satisfying \eqref{eq:curvature}, be even, and have
variance one. Independent copies are used for all scalar weights. The
\emph{strong-prior scales} are
\begin{equation}\label{eq:strong-scales}
 W^1_{ij}=\sigma\sqrt{N/d}\,\xi^1_{ij},\qquad
 W^s_{ij}=\sigma\xi^s_{ij}\ (2\leq s<L),\qquad
 W^L_j=\sigma N^{-1/2}\xi^L_j.
\end{equation}
The \emph{finite-scale exponent family} is
\begin{equation}\label{eq:gamma-scales}
 W^1_{ij}=N^{(1-\gamma)/2}\xi^1_{ij}/\sqrt d,\qquad
 W^s_{ij}=N^{-\gamma/2}\xi^s_{ij}\ (2\leq s<L),\qquad
 W^L_j=N^{-(\gamma+1)/2}\xi^L_j.
\end{equation}
These are exactly \eqref{eq:strong-scales} at $\sigma=N^{-\gamma/2}$.
At $\gamma=1$, the first-layer, internal, and readout variances are
$1/d$, $1/N$, and $1/N^2$.

\begin{theorem}[Strong-prior learning]\label{thm:strong-main}
Fix $L$ and a coordinate law as above. There is a polynomial $p_L$, with
coefficients allowed to depend on the law, such that the following holds.
Let $C=C_L(f)<\infty$, $N\geq C$, and $\varepsilon\in(0,1)$.
There is $\sigma_0>0$ such that for every $0<\sigma<\sigma_0$, every
input distribution $D$, every $\delta\in(0,1)$, and
\[
 n\geq\frac6\varepsilon
       \left[p_L(d,C)+\log\frac1\delta\right],
\]
with probability at least $1-\delta$,
$R_D(\Pi)\leq\varepsilon/3$ and $F_{2/3}(\Pi)\leq\varepsilon$.
The sample polynomial is independent of $N$. The threshold $\sigma_0$
may depend on the finite width, target, and desired accuracy; no uniform
small-scale threshold over all widths is asserted.
\end{theorem}

\begin{theorem}[Finite-scale learning for every $\gamma\geq1$]
\label{thm:finite-main}
Fix $L$, a coordinate law as above, and $\gamma\geq1$. There is a
polynomial $p_{L,\gamma}$ such that, for every $C=C_L(f)<\infty$,
every $N\geq C$, every distribution $D$ on $\mathcal X$, and every
$\varepsilon,\delta\in(0,1)$, the condition
\begin{equation}\label{eq:main-sample}
 n\geq\frac6\varepsilon\left[
 p_{L,\gamma}\!\left(d,C,\log\frac1\varepsilon\right)
 +\log\frac6{\varepsilon\delta}\right]
\end{equation}
implies, with probability at least $1-\delta$,
\[
 R_D(\Pi)\leq\varepsilon/3,\qquad F_{2/3}(\Pi)\leq\varepsilon.
\]
This is a finite-width statement for the actual non-Gaussian prior
\eqref{eq:gamma-scales}. The polynomial is independent of $N$, but its
degree and constants may depend on $L$, $\gamma$, and the coordinate law.
No uniformity as $\gamma\downarrow1$ is claimed for the $\gamma>1$ proof;
the endpoint has its own proof.
\end{theorem}

Theorems~\ref{thm:strong-main} and \ref{thm:finite-main} are noiseless,
statistical statements. They do not provide efficient posterior algorithms,
useful optimized numerical exponents, or polynomial dependence on depth.
A different fixed confidence threshold in $(0,1)$ follows from
\eqref{eq:confidence} by changing constants.

\section{Two elementary posterior estimates}\label{sec:posterior}
\begin{lemma}[A finite relevant class]\label{lem:finite-posterior}
Suppose a prior on finitely many functions assigns positive mass to $f$.
Let $\mathcal H$ contain $f$, with $\log|\mathcal H|\leq H$ and
$\pi(\mathcal H^c)/\pi(f)\leq\tau$. If
$n\geq a^{-1}[H+\log(1/\delta)]$, then its consistency posterior has
risk at most $a+\tau$ with sample probability at least $1-\delta$.
\end{lemma}
\begin{proof}
Each function of true error greater than $a$ survives all $n$ labels with
probability at most $e^{-na}$. A union bound excludes all such survivors
in $\mathcal H$ with the stated probability. For every sample the
posterior denominator is at least $\pi(f)$, so its mass outside
$\mathcal H$ is at most $\tau$.
\end{proof}

\begin{lemma}[Learning from target mass]\label{lem:target-posterior}
If $-\log\pi(f)\leq B$, then
$n\geq a^{-1}[B+\log(1/(\tau\delta))]$ implies posterior risk at most
$a+\tau$ with sample probability at least $1-\delta$.
\end{lemma}
\begin{proof}
Let $Z$ be posterior mass on functions of error greater than $a$. The
target lower-bounds the posterior denominator, whence
\[
 \E Z\leq\pi(f)^{-1}\sum_{g:\mathrm{err}_D(g)>a}
       \pi(g)(1-\mathrm{err}_D(g))^n\leq e^{B-na}.
\]
Markov's inequality gives $Z\leq\tau$ with probability at least
$1-\delta$. All other functions have error at most $a$.
\end{proof}

\section{A target-mass bound at every width}\label{sec:target-mass}
Throughout this section $\nu$ has the even density in
\eqref{eq:curvature}. Adding a constant to $V$ changes nothing, so take
$V(0)=0$. Evenness gives $V'(0)=0$ and hence
$m t^2/2\leq V(t)\leq Mt^2/2$.

\begin{lemma}[Concentration and translation]\label{lem:concentration}
A product of $\nu$ coordinates, conditioned on any centrally symmetric
convex set of positive probability, is centered and satisfies
\begin{equation}\label{eq:conditional-subgaussian}
 \E e^{\lambda u\cdot Z}\leq
        \exp\{\lambda^2\|u\|_2^2/(2m)\}.
\end{equation}
For the unconditioned product measure $\mu$, Euclidean enlargement obeys
\begin{equation}\label{eq:enlargement}
 \mu(A_r)\geq\Phi\bigl(\Phi^{-1}(\mu(A))+\sqrt m\,r\bigr).
\end{equation}
Finally, if $E$ is any centrally symmetric event of positive probability
and $a$ is deterministic, then
\begin{equation}\label{eq:translation-mass}
 \mu(E+a)\geq\mu(E)\exp\{-M\|a\|_2^2/2\}.
\end{equation}
\end{lemma}
\begin{proof}
Conditioning on a convex set preserves $m$-strong log-concavity. The
Brascamp--Lieb variance inequality \cite{brascamp-lieb}, applied after
any linear exponential tilt, bounds the Hessian of its log moment
generating function by $m^{-1}I$. Central symmetry makes the first
derivative at zero vanish. Integrating twice proves
\eqref{eq:conditional-subgaussian}. Convex indicator potentials follow
by approximation with smooth convex barriers.

For \eqref{eq:enlargement}, the Caffarelli contraction theorem
\cite{caffarelli} gives an $m^{-1/2}$-Lipschitz transport from standard
Gaussian measure to $\mu$. The inverse image of $A_r$ contains the
$\sqrt m\,r$ enlargement of the inverse image of $A$. Gaussian
isoperimetry \cite{ledoux} gives the displayed inequality.

For the last assertion let $Q$ be the law of $Z+a$ conditional on
$Z\in E$. Direct calculation of relative entropy gives
\[
 D(Q\Vert\mu)=-\log\mu(E)
 +\E\left[\sum_i(V(Z_i+a_i)-V(Z_i))\mid E\right].
\]
Taylor's theorem bounds the sum by
$\sum_i a_i V'(Z_i)+M\|a\|_2^2/2$. The conditional expectation of
its linear term is zero because $V'$ is odd and the conditional law is
invariant under simultaneous negation. Since $Q$ is supported on $E+a$,
$D(Q\Vert\mu)\geq-\log\mu(E+a)$. This proves
\eqref{eq:translation-mass}.
\end{proof}

\begin{lemma}[One-dimensional small balls]\label{lem:small-ball}
Every centered even log-concave variable $Y$ of variance at most $1/m$
satisfies
\[
 \Pr\{|Y|\leq\rho\}\geq c_m\rho\quad(0<\rho\leq1),
 \qquad c_m=\tfrac38\min\{1,\sqrt m\}.
\]
This applies to every normalized sum of disjoint coordinates of $Z$.
\end{lemma}
\begin{proof}
Let $R=2\max\{1,m^{-1/2}\}$. Chebyshev gives probability at least
$3/4$ in $[-R,R]$. Symmetric unimodality gives probability at least
$(\rho/R)(3/4)$ in $[-\rho,\rho]$. Log-concavity is preserved by
convolution and linear projection, and normalized independent sums retain
variance at most $1/m$.
\end{proof}

We first construct an exact clone of a target. Partition the $N$ neurons
of each hidden layer into $C$ groups of sizes $k_a\in\{\lfloor N/C\rfloor,
\lceil N/C\rceil\}$. Copy the corresponding teacher neuron throughout
each group. For an internal teacher weight $W_{ab}$ use $W_{ab}/k_b$
on every edge from group $b$ to group $a$. Divide the readout weight
from group $a$ by $k_a$. The cloned network computes exactly the same
function. Let $P$ project onto the group-constant subspace, with orthonormal
basis $u_a=k_a^{-1/2}\mathbf1_{\mathrm{group}\ a}$.

Define the weighted energy
\begin{equation}\label{eq:kappa}
 \kappa_N(w)=\frac dN\|W^1\|_F^2+
     \sum_{s=2}^{L-1}\|W^s\|_F^2+N\|W^L\|_2^2,
 \qquad \kappa_N(g)=\inf_{w\in U_{N,g}}\kappa_N(w).
\end{equation}
The first-layer clone contribution is at most $d(d+1)C^2$.
Because $\max k_a/\min k_a\leq2$ and $N/\min k_a\leq2C$,
its internal contribution is at most $2(L-2)C^4$ and its readout
contribution at most $2C^4$. Amplify the teacher readout by three; the
resulting exact clone has signed margin greater than $3/2$. A convenient
common bound for its energy is
\begin{equation}\label{eq:B0}
 B_0=d(d+1)C^2+2(L-2)C^4+36C^4.
\end{equation}
The larger final constant leaves harmless slack.

\begin{proposition}[Target mass]\label{prop:target-mass}
For each fixed $L$ there is a polynomially bounded function $A_L(d,C)$,
depending also on $m,M$, such that
\begin{equation}\label{eq:target-mass}
 \mu_{N,\gamma}(U_{N,f})\geq e^{-A_L(d,C)N^\gamma}
 \qquad(\gamma\geq1,\ N\geq C=C_L(f)).
\end{equation}
\end{proposition}
\begin{proof}
Put $t=N^{(1-\gamma)/2}\leq1$. In standardized coordinates impose the
centered cylinder $\mathcal C_\rho$ consisting of the following constraints:
all first-layer and readout coordinates have absolute value at most $\rho$;
each internal row obeys $|\langle Z_{j,*},u_a\rangle|\leq\rho$ for
every group $a$. These constraints involve independent random variables:
within an internal row the sums concern disjoint groups. There are
\[
 r_0=N[d+2+(L-2)C]
\]
constraints. Lemma~\ref{lem:small-ball} gives
$\mu(\mathcal C_\rho)\geq(c_m\rho)^{r_0}$.

Conditional on the cylinder, every internal matrix is centered and
$m$-strongly log-concave. For unit vectors $u,v$, the scalar $u^TZv$
is therefore $m^{-1/2}$-sub-Gaussian. Two $1/4$-nets, each with at most
$9^N$ elements, give
\[
 \Pr\{\|Z\|_{\mathrm{op}}>2q\sqrt N\mid\mathcal C_\rho\}
 \leq2\exp\{(2\log9-mq^2/2)N\}.
\]
Take
\[
 b_{m,L}=2\sqrt{\frac2m(2\log9+\log(4L))}.
\]
Each conditional failure probability for
$\|Z\|_{\mathrm{op}}\leq b_{m,L}\sqrt N$ is at most $1/(2L)$.
Let $E$ be the cylinder intersected with these bounds for all internal
layers. It is centrally symmetric and has
$\mu(E)\geq\tfrac12(c_m\rho)^{r_0}$.

Translate this event by the amplified clone in standardized coordinates.
The internal weight perturbation decomposes as
\[
 \frac t{\sqrt N}Z=R+E_\perp,
 \qquad R=\frac t{\sqrt N}ZP,
 \qquad E_\perp=\frac t{\sqrt N}Z(I-P),
\]
where
$\|R\|_{\mathrm{op}}\leq\rho t\sqrt C$,
$E_\perp P=0$, and $\|E_\perp\|_{\mathrm{op}}\leq b_{m,L}t$.
We detail stability to make the role of the cylinder explicit. Write
$e_s=N^{-1/2}\|h^s-h_*^s\|_2$, with $h_*^s$ the exact clone's
activations. Its internal matrices have norm at most $\sqrt2C^2$, and
$h_*^s\in\mathrm{range}(P)$. The first layer has
\[
 e_1\leq \rho t(d+1)/\sqrt d\leq\rho(d+1).
\]
For $s\geq2$, boundedness and the contraction property of tanh give
\[
 e_s\leq(\sqrt2C^2+b_{m,L})e_{s-1}+\rho\sqrt C.
\]
Set $A_0=2+3C^2+b_{m,L}$. Induction bounds
$e_{L-1}\leq\rho(d+L+1)A_0^{L-2}$.
The cloned readout has $\sqrt N\|v_*\|_2\leq3\sqrt2 C^2$,
and the perturbation readout has $\ell^1$ norm at most $\rho t\leq\rho$.
Consequently the output error is at most
\[
 3\sqrt2 C^2\rho(d+L+1)A_0^{L-2}+\rho.
\]
Choose
\[
 \rho=[2+9C^2(d+L+1)A_0^{L-2}]^{-1}.
\]
The last error is strictly smaller than $1/2$, so the translated event
lies in $U_{N,f}$, indeed with signed margin greater than one.

If $a$ is the standardized translation, then
$\|a\|_2^2=N^\gamma\kappa_N(w_*)\leq B_0N^\gamma$.
Lemma~\ref{lem:concentration} gives
\[
 -\log\mu_{N,\gamma}(U_{N,f})
 \leq\tfrac M2B_0N^\gamma
 +N[d+2+(L-2)C]\log\frac1{c_m\rho}+\log2.
\]
Since $N\leq N^\gamma$, one may take
\begin{equation}\label{eq:Aexplicit}
 A_L(d,C)=1+\tfrac M2B_0+
 [d+2+(L-2)C]\log\frac1{c_m\rho}+\log2.
\end{equation}
At fixed depth and fixed curvature constants this has a polynomial
majorant in $d,C$.
\end{proof}

\section{The strong-prior limit and quadratic-tail envelopes}\label{sec:strong}
We first record the deterministic counting consequence of
Appendix~\ref{sec:norm-appendix}.
\begin{lemma}[A width-independent energy count]\label{lem:energy-count}
For every fixed $L$ there is a polynomially bounded $H_L(d,K)$ such that
\[
 \log\#\{g:\kappa_N(g)\leq K\}\leq H_L(d,K)
 \qquad(K\geq1),
\]
uniformly over all widths. The same bound holds for the union of these
function classes over widths.
\end{lemma}
\begin{proof}
Take a representative of energy less than $K+1$ for each function. The
compression theorem in Appendix~\ref{sec:norm-appendix} provides an
implementation of width $b=m_L(d,K+1)$, strict margin $1/2$, and weights
at most $b$, where $b$ is polynomially bounded at fixed depth. Round each
weight to a multiple of $(8\Lambda)^{-1}$, where
\[
 P=b(d+2)+(L-2)b^2,\qquad
 \Lambda=(d+Lb)[1+b(b+1)]^{L-1}.
\]
A layerwise maximum-norm induction bounds the output change by $1/8$;
this is also proved in Lemma~\ref{lem:tilt-count}. Thus all signs survive.
There are at most $(3+16b\Lambda)^P$ rounded arrays. Take
$H_L(d,K)=P\log(3+16b\Lambda)$. This counts one common compressed family
for every original width.
\end{proof}

The next statement permits many priors outside the curvature family.
\begin{proposition}[Two-sided quadratic envelopes]\label{prop:envelopes}
In the strong-prior scales \eqref{eq:strong-scales}, replace standardized
coordinates by independent densities $q_j$ satisfying
\begin{equation}\label{eq:envelopes}
 c_j e^{-b u^2}\leq q_j(u)\leq C_j e^{-a u^2}
 \quad(u\in\R),\qquad 0<a\leq b<\infty,
\end{equation}
where $c_j,C_j>0$ are finite. Assume $b/a$ is bounded independently of
width and coordinate. Then Theorem~\ref{thm:strong-main} holds, with a
sample polynomial depending also on $b/a$. The prefactors $c_j,C_j$ need
not be uniform in width; they affect only the fixed-width small-scale
threshold.
\end{proposition}
\begin{proof}
At fixed width there are finitely many coordinates. If $s_j$ denotes
$\sqrt{N/d}$, $1$, or $N^{-1/2}$ according to layer, then
$\kappa_N(w)=\sum_j(w_j/s_j)^2$. A strict target witness obtained by
cloning has energy at most $B_0$ from \eqref{eq:B0}. For every
$\zeta>0$, a sufficiently small positive-volume neighborhood of the
witness lies in $U_{N,f}$ and has energy at most $B_0+\zeta$.
Integrating the lower density envelope there gives
\begin{equation}\label{eq:strong-lower}
 \liminf_{\sigma\downarrow0}\sigma^2
       \log\mu_{N,\sigma}(U_{N,f})\geq-bB_0.
\end{equation}
All powers of $\sigma$ and fixed density prefactors disappear after
multiplication by $\sigma^2$.

For any $g$ and $\theta\in(0,1)$, the upper envelope and the inequality
$\kappa_N(w)\geq\kappa_N(g)$ on $U_{N,g}$ give
\[
 \mu_{N,\sigma}(U_{N,g})
 \leq C_{N,\theta}
  \exp\{-(1-\theta)a\kappa_N(g)/\sigma^2\}.
\]
Here the remaining integral of
$\sigma^{-P}\exp\{-\theta a\kappa_N(w)/\sigma^2\}$ is a constant
independent of $\sigma$. Letting $\theta\downarrow0$ proves
\begin{equation}\label{eq:strong-upper}
 \limsup_{\sigma\downarrow0}\sigma^2
       \log\mu_{N,\sigma}(U_{N,g})\leq-a\kappa_N(g).
\end{equation}
Set $K=(b/a)B_0+1$ and $\mathcal H=\{g:\kappa_N(g)\leq K\}$.
Each function outside $\mathcal H$ has mass negligible relative to $f$
by \eqref{eq:strong-lower}--\eqref{eq:strong-upper}. There are only
finitely many Boolean functions, so their total ratio tends to zero.
Global-margin conditioning cancels in this ratio. Choose $\sigma_0$ so
that it is at most $\varepsilon/6$. Lemma~\ref{lem:energy-count} and
Lemma~\ref{lem:finite-posterior}, with $a=\tau=\varepsilon/6$, give
Gibbs risk at most $\varepsilon/3$ as soon as
\[
 n\geq\frac6\varepsilon[H_L(d,K)+\log(1/\delta)].
\]
The threshold is independent of $D,n,\delta$. Finally use
\eqref{eq:confidence}. For the curvature family, integrating the bounds
on $V''$ gives \eqref{eq:envelopes} with $a=m/2$, $b=M/2$;
this proves Theorem~\ref{thm:strong-main}.
\end{proof}

\begin{proposition}[Unchanged strong-prior ordering]\label{prop:ordering}
Suppose the finitely many standardized coordinate densities are positive
and continuous, and for a common $a>0$ satisfy
\[
 \lim_{u\to+\infty}\frac{-\log q_j(u)}{u^2}
 =\lim_{u\to-\infty}\frac{-\log q_j(u)}{u^2}=a.
\]
Then for every realizable function $g$,
\[
 \lim_{\sigma\downarrow0}\sigma^2
     \log\mu_{N,\sigma}(U_{N,g})=-a\kappa_N(g).
\]
Thus its leading strong-prior ordering is exactly the weighted quadratic
energy ordering. That energy is polynomially equivalent to $C_L(g)$
whenever $N\geq C_L(g)$.
\end{proposition}
\begin{proof}
For any $\zeta\in(0,a)$ the tail asymptotics, positivity, and continuity
supply global upper and lower envelopes with exponents $a-\zeta$ and
$a+\zeta$. Apply the preceding integral estimates to witnesses arbitrarily
close to the infimum $\kappa_N(g)$, then let $\zeta\downarrow0$.
Cloning bounds energy polynomially in $d,C_L(g)$; deterministic
compression bounds $C_L(g)$ polynomially in $d,\kappa_N(g)+1$.
\end{proof}

Examples of \eqref{eq:envelopes} include Gaussian densities multiplied by
positive bounded factors, and centered Gaussian scale mixtures with scales
in a fixed compact subinterval of $(0,\infty)$. Neither symmetry nor
log-concavity is needed for this strong-prior proposition. The density
$\exp\{-u^2/2-b\log\cosh u\}$ also satisfies
Proposition~\ref{prop:ordering}, because its additional log density is only
linear in $|u|$ at infinity. Variance standardization changes the common
quadratic coefficient but not the ordering.

\section{Finite-scale learning for \texorpdfstring{$\gamma>1$}{gamma > 1}}\label{sec:above}
\begin{lemma}[A non-Gaussian singular-value tail]\label{lem:spectral}
Let $B=\sigma Z$ be an $N\times N$ matrix with independent $\nu$
coordinates. If $a\geq16\sigma\sqrt{N/m}$ and $B_{>a}$ retains
exactly the singular values greater than $a$, then
\[
 \Pr\{\|B_{>a}\|_F^2>K\}\leq e^{-mK/(8\sigma^2)}.
\]
\end{lemma}
\begin{proof}
The net proof in Section~\ref{sec:target-mass} gives
\[
 \Pr\{\|B\|_{\mathrm{op}}>8\sigma\sqrt{N/m}\}
 \leq2e^{(2\log9-8)N}<1/2.
\]
The convex ball $\mathcal B=\{T:\|T\|_{\mathrm{op}}\leq a/2\}$
therefore has probability at least $1/2$. Its Frobenius distance from $B$
is
\[
 D(B)^2=\sum_j(s_j(B)-a/2)_+^2.
\]
The enlargement inequality \eqref{eq:enlargement}, after scaling by
$\sigma$, implies $\Pr\{D(B)>r\}\leq e^{-mr^2/(2\sigma^2)}$.
Since $s>a$ implies $s\leq2(s-a/2)$,
$\|B_{>a}\|_F\leq2D(B)$. Take $r=\sqrt K/2$.
\end{proof}

For later use, \eqref{eq:conditional-subgaussian} and a one-dimensional
Gaussian integration identity give
\begin{equation}\label{eq:square-moment}
 \E e^{m\xi^2/4}\leq\sqrt2.
\end{equation}
Indeed $e^{m x^2/4}=\E_G e^{\sqrt{m/2}Gx}$, with $G$ standard normal;
condition first on $G$. For $D_0$ independent coordinates, Markov's
inequality consequently gives
\begin{equation}\label{eq:sum-squares}
 \Pr\{\textstyle\sum_{j=1}^{D_0}\xi_j^2>u\}
 \leq2^{D_0/2}e^{-mu/4}\leq e^{-mu/8}
 \quad(mu\geq8D_0).
\end{equation}

\begin{lemma}[Aggregate structural tail]\label{lem:structural}
Set
\[
 a_K=[8L(\sqrt K+1)^{L-2}]^{-1},\qquad J_K=(L+3)K,
 \qquad K\geq1.
\]
If
\begin{equation}\label{eq:width-above}
 N^{\gamma-1}\geq8(d+1)/m,\qquad
 N^{\gamma-1}\geq256/(ma_K^2),
\end{equation}
then
\[
 \mu_{N,\gamma}\!\left(\bigcup_{g:\kappa_N(g)>J_K}U_{N,g}\right)
 \leq L e^{-mKN^\gamma/8}.
\]
\end{lemma}
\begin{proof}
Consider the event
\[
 \tfrac dN\|W^1\|_F^2\leq K,\quad
 N\|W^L\|_2^2\leq K,\quad
 \|(W^s)_{>a_K}\|_F^2\leq K\ (2\leq s<L).
\]
The first and last weighted norms are $N^{-\gamma}$ times sums of
$N(d+1)$ and $N$ standardized squares. Equation~\eqref{eq:sum-squares}
and the first width condition bound their tails by $e^{-mKN^\gamma/8}$.
The second width condition permits Lemma~\ref{lem:spectral} at each
internal layer. A union bound bounds failure of the event by
$Le^{-mKN^\gamma/8}$.

For a draw in this event, retain only internal singular values greater
than $a_K$. The retained matrices have operator norm at most $\sqrt K$;
each discarded matrix has norm at most $a_K$. If $e_s$ is the Euclidean
activation error, then $e_1=0$ and
$e_s\leq\sqrt K e_{s-1}+a_K\sqrt N$. The readout norm is at most
$\sqrt{K/N}$, so the uniform output error is at most
\[
 a_K\sum_{j=1}^{L-2}K^{j/2}\leq1/8.
\]
For $L=2$ there is no change. A strict-margin draw therefore retains
margin greater than $3/8$. Doubling its readout restores margin greater
than $1/2$ and gives energy at most
$K+(L-2)K+4K=J_K$. Every full-margin draw outside the asserted function
class must belong to the exceptional parameter event, proving the lemma.
\end{proof}

\begin{proof}[Proof of Theorem~\ref{thm:finite-main} for $\gamma>1$]
Write $A=A_L(d,C)$ and choose
\[
 K=\max\{1,\tfrac8m[A+\log(6L/\varepsilon)]\},\qquad
 T=\left\lceil\max\{C,[8(d+1)/m]^{1/(\gamma-1)},
                [256/(ma_K^2)]^{1/(\gamma-1)}\}\right\rceil.
\]
For $N\geq T$, divide the aggregate tail by the target lower bound
\eqref{eq:target-mass}. Global-margin normalization cancels and gives
\[
 \frac{\pi(\{g:\kappa_N(g)>J_K\})}{\pi(f)}
 \leq L e^{-N^\gamma(mK/8-A)}\leq\varepsilon/6.
\]
Necessarily $f$ belongs to the remaining class, since otherwise the ratio
would be at least one. Its log cardinality is at most $H_L(d,J_K)$.
Lemma~\ref{lem:finite-posterior}, with $a=\tau=\varepsilon/6$, gives
the claimed risk for
$n\geq(6/\varepsilon)[H_L(d,J_K)+\log(1/\delta)]$.

For $C\leq N<T$, the target mass gives
$-\log\pi(f)\leq AN^\gamma\leq AT^\gamma$.
Lemma~\ref{lem:target-posterior} gives the same risk if
$n\geq(6/\varepsilon)[AT^\gamma+\log(6/(\varepsilon\delta))]$.
At fixed $L,\gamma,m,M$, all of $K,T,T^\gamma$ have polynomial
majorants in $d,C,\log(1/\varepsilon)$. Majorize
$\max\{H_L(d,J_K),AT^\gamma\}$ by $p_{L,\gamma}$ and use
\eqref{eq:confidence}. This covers every $N\geq C$.
\end{proof}

\section{The mean-field endpoint \texorpdfstring{$\gamma=1$}{gamma = 1}}\label{sec:meanfield}
At this endpoint the internal matrices cannot be discarded by the preceding
spectral argument. We instead view one neuron's values over all Boolean
inputs as a random vector, called its field, and track covariances of these
vectors through the layers. Conditional on the preceding layer, a
non-Gaussian row and its Gaussian comparison have the same covariance
but different distributions. The following estimate compares their log
partition functions uniformly in the field's dimension, which will be
$2^d$.

For matrices, $\|\Gamma\|_{1,\mathrm{ent}}$ denotes the sum of absolute
entries, and $\|Q\|_{\max}$ denotes the largest absolute entry. The
first norm controls tilt parameters and the second measures covariance
error; their pairing avoids a factor equal to the number of input pairs.

\begin{lemma}[A relative Lindeberg estimate]\label{lem:lindeberg}
Let $a_1,\ldots,a_N\in\R^s$ obey $\|a_j\|_\infty\leq N^{-1/2}$.
Let $\xi_j$ be independent centered variance-one variables with law $\nu$,
and let $g_j$ be independent standard Gaussians. Define
\[
 Z_\nu=\sum_j a_j\xi_j,\qquad Z_G=\sum_j a_jg_j,\qquad
 F_\Gamma(z)=\tanh(z)^T\Gamma\tanh(z),
\]
where $\|\Gamma\|_{1,\mathrm{ent}}\leq G$. Write
\[
 P(G)=24G+36G^2+8G^3,\qquad
 J_\nu=\E[|\xi|^3e^{|\xi|}]+\E[|g|^3e^{|g|}],\qquad
 c_\nu=eJ_\nu/6.
\]
If $N\geq4G^2$, then
\begin{equation}\label{eq:lindeberg}
 \left|\log\E e^{F_\Gamma(Z_\nu)}-
             \log\E e^{F_\Gamma(Z_G)}\right|
 \leq c_\nu P(G)/\sqrt N.
\end{equation}
The constant is finite for the curvature family. The conclusion requires
only the displayed moment, centering, and variance assumptions.
\end{lemma}
\begin{proof}
Using $|\tanh|,|\tanh'|\leq1$, $|\tanh''|\leq2$, and
$|\tanh'''|\leq6$, differentiation in a direction $a$ with
$\|a\|_\infty\leq N^{-1/2}$ gives
\[
 |D_aF_\Gamma|\leq2G/\sqrt N,\quad
 |D_a^2F_\Gamma|\leq6G/N,\quad
 |D_a^3F_\Gamma|\leq24G/N^{3/2}.
\]
Thus $|D_a^3e^{F_\Gamma}|\leq P(G)N^{-3/2}e^{F_\Gamma}$.
Fix the residual vector $z$ in a one-coordinate replacement and put
$f_z(t)=e^{F_\Gamma(z+at)}$ and $\lambda=2G/\sqrt N\leq1$.
The first derivative bound implies
\[
 e^{-\lambda|t|}f_z(0)\leq f_z(t)\leq e^{\lambda|t|}f_z(0).
\]
Taylor expansion through order two with integral remainder yields
\[
 |f_z(t)-f_z(0)-tf_z'(0)-t^2f_z''(0)/2|
 \leq\frac{P(G)}{6N^{3/2}}|t|^3e^{\lambda|t|}f_z(0).
\]
The first three averaged terms match for $\xi$ and $g$. The difference
of the one-coordinate expectations is therefore at most
$P(G)J_\nu f_z(0)/(6N^{3/2})$. Both expectations are at least
$e^{-1}f_z(0)$, since Jensen and variance one give
$\E e^{-\lambda|\xi|}\geq e^{-\lambda\E|\xi|}\geq e^{-1}$.
The same holds for $g$.

Averaging over any independent random residual $z$ preserves both bounds.
The inequality $|\log u-\log v|\leq|u-v|/\min\{u,v\}$ shows that
the logarithms differ by at most $c_\nu P(G)N^{-3/2}$ at this
replacement. Telescoping over all $N$ coordinates proves
\eqref{eq:lindeberg}. No inverse global partition-function bound was
used; such a bound would have introduced an unnecessary exponential in
$G$.
\end{proof}

Let $s=2^d$, $r=L-1$, and index fields and covariance matrices by all
inputs. Write
\[
 Q_{\ell,N}=\frac1N\sum_{i=1}^N h_i^\ell(h_i^\ell)^T.
\]
For a law $\mu$ on $[-1,1]^s$, define
\begin{equation}\label{eq:certificate}
 \mathcal J_G(\widehat Q;\mu)=
 \max_{\|\Gamma\|_{1,\mathrm{ent}}\leq G}
 \left\{\langle\Gamma,\widehat Q\rangle
       -\log\E_{h\sim\mu}e^{h^T\Gamma h}\right\}.
\end{equation}
Appendix~\ref{sec:sparse-dual} proves, for any law $\mu$, that an
empirical second moment from $N$ independent rows satisfies
\begin{equation}\label{eq:certificate-tail}
 \Pr\{\mathcal J_G(\widehat Q;\mu)>K\}
 \leq\exp\{q\log(2s^2+1)-N(K-1)\},\qquad q=\lceil G^2/2\rceil.
\end{equation}
It also proves that $\mathcal J_G\leq K$ supplies a quadratic
exponential tilt, of parameter norm at most $G$, whose second moment is
within $K/G$ of $\widehat Q$ in entrywise maximum norm. These facts
are distribution-free.

Let $\mu_0$ be the actual first-layer activation-field law, obtained
from independent $\nu$ weights divided by $\sqrt d$. Keep this law
exactly; its dimension need not tend to infinity. For later layers write
$\mathsf M_Q$ for the activation law of a Gaussian field of covariance
$Q$, and $\mathsf F_\Gamma(Q)$ for its second moment after tilting by
$e^{h^T\Gamma h}$. Lemma~\ref{lem:tilted-covariance} proves
\begin{equation}\label{eq:covariance-stability}
 \|\mathsf F_\Gamma(Q)-\mathsf F_\Gamma(Q')\|_{\max}
 \leq8(1+G)^2\|Q-Q'\|_{\max}.
\end{equation}
This comparison concerns an auxiliary Gaussian family; the Bayesian
learner still has non-Gaussian weights at every layer.

Conditional on the preceding activations $h_1,\ldots,h_N$, an internal
row has preactivation $N^{-1/2}\sum_j\xi_jh_j$. Apply
Lemma~\ref{lem:lindeberg} with $a_j=h_j/\sqrt N$. Its Gaussian
comparison has covariance exactly $Q=N^{-1}\sum_jh_jh_j^T$.
If $\mu_{\mathrm{act}}$ is the true conditional activation law, then
\begin{equation}\label{eq:certificate-transfer}
 \mathcal J_G(\widehat Q;\mathsf M_Q)
 \leq\mathcal J_G(\widehat Q;\mu_{\mathrm{act}})
       +c_\nu P(G)/\sqrt N.
\end{equation}
This step does not assume that the actual row law is determined by $Q$.

\subsection{A likely finite class}
The covariance comparisons now produce a finite class containing every
global-margin function outside a small parameter event. We choose the
accuracy budgets backward through depth, so that errors accumulated by
the forward reconstruction stay below the output margin. These budgets
remain polynomial only when depth is fixed.
Let $K\geq4$ and $\bar K=K+1$. Choose the budgets backward:
\begin{equation}\label{eq:budgets}
 G_r=32r\bar K^2,\qquad
 G_\ell=8G_{\ell+1}(1+G_{\ell+1})^2\quad(\ell<r),
 \qquad b_\ell=8(1+G_\ell)^2.
\end{equation}
Set $q_1=\lceil G_1^2/2\rceil$, $\ell_0=\log(2s^2+1)$, and
\begin{equation}\label{eq:meanfield-threshold}
 T(K)=\left\lceil\max\left\{1,\frac{2q_1\ell_0}{K-1},
                    4G_1^2,c_\nu^2P(G_1)^2\right\}\right\rceil.
\end{equation}
These quantities have polynomial majorants in $d,K$ at each fixed depth:
$\ell_0\leq(2d+2)\log2$ and the recursion in \eqref{eq:budgets} has
only the fixed number $r$ of steps.

\begin{lemma}[Likely proxy covariances]\label{lem:proxy}
There is a constant $c>0$, depending only on the coordinate law, such
that for all sufficiently large fixed-law $K$ and $N\geq T(K)$,
except on a parameter event of probability at most $Le^{-cNK}$, the
following hold simultaneously: $N\|W^L\|_2^2\leq K$, and there are
symmetric matrices $\Gamma_\ell$ with
$\|\Gamma_\ell\|_{1,\mathrm{ent}}\leq G_\ell$ such that the tower
\[
 Q_1^*=\frac{\E_{h\sim\mu_0}[hh^Te^{h^T\Gamma_1h}]}
                    {\E_{h\sim\mu_0}e^{h^T\Gamma_1h}},
 \qquad Q_\ell^*=\mathsf F_{\Gamma_\ell}(Q_{\ell-1}^*)\ (\ell\geq2)
\]
satisfies $\|Q_{r,N}-Q_r^*\|_{\max}\leq1/(32\bar K)$.
\end{lemma}
\begin{proof}
Apply \eqref{eq:certificate-tail} to each actual conditional row law,
including $\mu_0$ in the first layer. The threshold gives each failure
probability at most
$e^{-N(K-1)/2}\leq e^{-NK/4}$. By
\eqref{eq:certificate-transfer} and \eqref{eq:meanfield-threshold},
every internal Gaussian certificate is at most $K+1=\bar K$.
The first-layer certificate is at most $K$ and needs no comparison.
The bounded-dual lift in Appendix~\ref{sec:sparse-dual} therefore
supplies the tilt matrices with local covariance error at most
$\bar K/G_\ell$.

Let $e_\ell=\|Q_{\ell,N}-Q_\ell^*\|_{\max}$.
Then $e_1\leq\bar K/G_1$, and
$e_\ell\leq\bar K/G_\ell+b_\ell e_{\ell-1}$ for $\ell\geq2$,
by \eqref{eq:covariance-stability}. The backward budgets give
\[
 e_r\leq\sum_{\ell=1}^r\frac{\bar K}{G_\ell}
                  \prod_{j=\ell+1}^r b_j
       =\frac1{32\bar K}.
\]
For the readout, \eqref{eq:square-moment} gives
\[
 \Pr\{N\|W^L\|_2^2>K\}
 \leq\exp\{-mNK/4+(N/2)\log2\}\leq e^{-mNK/8}
\]
when $K\geq4\log2/m$. Thus $c=\min\{1/4,m/8\}$ works after
requiring $K\geq\max\{4,4\log2/m\}$.
\end{proof}

\begin{lemma}[A bounded dual readout]\label{lem:readout-dual}
Suppose $H$ is an $s\times N$ feature matrix and a readout $v$ obeys
$g_i(Hv)_i>1/2$ and $N\|v\|_2^2\leq K$. Put $Q=HH^T/N$.
There is $\beta\in\R^s$ with $\|\beta\|_1\leq2K$ and
$g_i(Q\beta)_i\geq1/2$ for every $i$.
\end{lemma}
\begin{proof}
Minimize $N\|u\|_2^2/2$ subject to $g_i(Hu)_i\geq1/2$.
The strictly feasible $v$ gives the convex optimality conditions with
multipliers $\lambda_i\geq0$:
$Nu=H^T\beta$, $\beta_i=g_i\lambda_i$, and
$\lambda_i[g_i(Hu)_i-1/2]=0$. Hence $Hu=Q\beta$ and
\[
 \tfrac12\|\beta\|_1=\beta^THu=N\|u\|_2^2\leq K.
\]
No inversion or nonsingularity of $Q$ is needed.
\end{proof}

\begin{proposition}[Mean-field aggregate mass bound]\label{prop:mf-aggregate}
For each $K$ as in Lemma~\ref{lem:proxy}, there is a finite class
$\mathcal H_K$, independent of learner width, with
$\log|\mathcal H_K|\leq\widetilde H_{L,\nu}(d,K)$ for a
polynomially bounded function at fixed depth, such that
\[
 \mu_{N,1}\!\left(\bigcup_{g\notin\mathcal H_K}U_{N,g}\right)
 \leq Le^{-cNK}\qquad(N\geq T(K)).
\]
\end{proposition}
\begin{proof}
Fix a full-margin draw in the likely set of Lemma~\ref{lem:proxy}.
Lemma~\ref{lem:readout-dual} gives $\|\beta\|_1\leq2K$ and margin
at least $1/2$ for $Q_{r,N}\beta$. Replacing $Q_{r,N}$ by $Q_r^*$
changes the output by at most $2K/(32\bar K)\leq1/16$.
Consequently $2Q_r^*\beta$ has signed margin at least $7/8$.

Proposition~\ref{prop:tilt-compress} realizes this tower within $1/8$
by a finite depth-$L$ network with width and weights polynomially bounded
in $d,\bar K,G_1$. Its first hidden layer is sampled from the true
first-layer law with the tilt; subsequent layers use Gaussian rows to
realize the auxiliary Gaussian tower. This constructs a deterministic
short description of the sign function, without changing the learner.
The realized margin is at least $3/4$. Define $\mathcal H_K$ to be all
sign functions admitting these common width and weight bounds with margin
at least $3/4$, padding by zero neurons to the common width when needed.
Lemma~\ref{lem:tilt-count} bounds this class's logarithmic cardinality
polynomially in $d,K$ at fixed depth. In particular every member of
$\mathcal H_K$ has the stated bounded implementation; the grid is used
only to count such functions. Every full-margin draw outside this class
must fail the likely event, giving the stated aggregate bound.
\end{proof}

\begin{proof}[Proof of Theorem~\ref{thm:finite-main} for $\gamma=1$]
Let $A=A_L(d,C)$ and choose
\[
 K\geq\max\{4,4\log2/m,\ c^{-1}[A+\log(6L/\varepsilon)]\}.
\]
For $N\geq T(K)$, Proposition~\ref{prop:mf-aggregate} and the target
mass give
\[
 \frac{\pi(\mathcal H_K^c)}{\pi(f)}
 \leq Le^{-N(cK-A)}\leq\varepsilon/6.
\]
The target therefore belongs to $\mathcal H_K$.
Lemma~\ref{lem:finite-posterior} gives risk at most $\varepsilon/3$
provided $n\geq(6/\varepsilon)[\widetilde H_{L,\nu}(d,K)+
\log(1/\delta)]$.
For every remaining width $C\leq N<T(K)$, the target lower bound gives
$-\log\pi(f)\leq AT(K)$. Lemma~\ref{lem:target-posterior} supplies
the same risk at sample size
$(6/\varepsilon)[AT(K)+\log(6/(\varepsilon\delta))]$.
Both $\widetilde H_{L,\nu}(d,K)$ and $AT(K)$ have polynomial
majorants in $d,C,\log(1/\varepsilon)$ at fixed depth. Taking a common
majorant and applying \eqref{eq:confidence} completes the proof.
\end{proof}

\section{Prior mass and circuit complexity}\label{sec:prior-complexity}
The preceding arguments also compare function probability with bounded-weight
implementation complexity. This is stronger than merely having a sample
bound, although it does not specify the relative order of functions whose
complexities are close.

\begin{corollary}[A polynomial prior-complexity comparison]
\label{cor:prior-complexity}
Fix $L$, $\gamma\geq1$, and a coordinate law in the main family. For a
function with $\mu_{N,\gamma}(U_{N,f})>0$, define its raw mass cost
\[
 b_{N,\gamma}(f)=-N^{-\gamma}\log\mu_{N,\gamma}(U_{N,f}).
\]
There is a polynomial $P_{L,\gamma}$, independent of width, such that
\[
 C_L(f)\leq P_{L,\gamma}(d,1+b_{N,\gamma}(f)).
\]
Conversely, whenever $N\geq C_L(f)$,
\[
 b_{N,\gamma}(f)\leq A_L(d,C_L(f)).
\]
The two functions of complexity on the right have polynomial majorants at
fixed depth and exponent. Global-margin normalization is common to all
functions, so using raw masses does not change their ordering at a fixed
width and exponent.
\end{corollary}
\begin{proof}
The converse is Proposition~\ref{prop:target-mass}. Put
$b=b_{N,\gamma}(f)\geq0$. If $\gamma>1$, choose
$K\geq\max\{1,(8/m)[b+\log(2L)]\}$. Above the width threshold
in \eqref{eq:width-above}, Lemma~\ref{lem:structural} makes the total
mass outside the bounded-energy class at most
\[
 Le^{-[b+\log(2L)]N^\gamma}<e^{-bN^\gamma}
       =\mu_{N,\gamma}(U_{N,f}).
\]
Thus $U_{N,f}$ contains a good draw and $\kappa_N(f)\leq J_K$.
Deterministic compression bounds $C_L(f)$ polynomially in $d,1+b$.
The threshold itself is polynomially bounded in these variables for fixed
$L,\gamma,m$.

At $\gamma=1$, choose $K$ above the fixed-law thresholds with
$cK\geq b+\log(2L)$. For $N\geq T(K)$, the same mass comparison in
Proposition~\ref{prop:mf-aggregate} implies $f\in\mathcal H_K$.
The constructive proof of that proposition gives a network of width and
weights polynomially bounded in $d,K$. Padding by zero neurons if
necessary shows that $C_L(f)$ has the required bound.

It remains to handle widths below the respective polynomial thresholds.
Let $D_N=N(d+2)+(L-2)N^2$ be the number of scalar parameters, and let
$S$ be the sum of their standardized squares. Exactly,
$S=N^\gamma\kappa_N(w)$. By the first inequality in
\eqref{eq:sum-squares}, at
\[
 u=\frac{4bN^\gamma+2D_N\log2+4}{m}
\]
one has $\Pr\{S>u\}\leq e^{-bN^\gamma-1}$.
This is strictly less than the mass of $U_{N,f}$, so that region contains
a point with
\[
 \kappa_N(w)\leq\frac u{N^\gamma}
 \leq\frac{4b+2D_N\log2+4}{m}.
\]
Since $N$ is below a polynomial threshold in $d,1+b$, this is a
polynomial energy bound. Deterministic compression proves the assertion
in the remaining width range.
\end{proof}

This comparison concerns the stated bounded-weight neural-network
complexity at fixed depth. It does not identify prior probability with an
exponential of ordinary Boolean circuit description length, and it does
not assert that posterior inference returns a minimum circuit. The exact
strong-prior rate under quadratic recession, rather than merely its
polynomial comparison, is Proposition~\ref{prop:ordering}.


\section{Scope and remaining questions}
The assumptions identify a broad sufficient family rather than a
classification of non-Gaussian priors. In particular, Laplace, Student,
uniform, and pure quartic-exponential coordinate laws are not covered by
Theorem~\ref{thm:finite-main}. Strong log-concavity alone would allow
superquadratic tails, whereas the upper curvature bound here prices a
translated target quadratically. A compactly supported coordinate law has
a more immediate difficulty in the fixed-width strong-prior limit: for
sufficiently small scale all possible outputs may fail the global margin,
leaving the conditioned learner undefined.

The strong-prior theorem actually needs only the two-sided Gaussian
envelopes \eqref{eq:envelopes} and permits asymmetric and non-log-concave
densities; see
Section~\ref{sec:strong}. At finite scale, symmetry supplies centering in
the translated conditional law, while strong log-concavity controls that
law even after a convex cylinder restriction. Removing either condition
requires a replacement for that target-mass argument. The mean-field
comparison itself needs substantially less: centered variance-one internal
coordinates with $\E[|\xi|^3e^{|\xi|}]<\infty$, suitable square-exponential
moments at the boundary layers, and a separate target-mass lower bound.

The theorems preserve the original global margin restriction. They do not
assert the same result after conditioning only on margins at training
points, under label noise, or for ReLU activations. Neither do they settle
which non-Gaussian laws at their natural width scales retain exactly the
same ordering of functions by prior probability, beyond the polynomial
comparison proved in Corollary~\ref{cor:prior-complexity}. For laws with asymptotically
quadratic log density, the strong-prior ordering is exactly the same
weighted quadratic-energy ordering; see Proposition~\ref{prop:ordering}.

\begingroup
\small
\begin{thebibliography}{9}
\setlength{\itemsep}{2pt}
\bibitem{hoeffding} W. Hoeffding.
Probability inequalities for sums of bounded random variables.
\emph{Journal of the American Statistical Association} 58 (1963), 13--30.
\bibitem{brascamp-lieb} H. J. Brascamp and E. H. Lieb.
On extensions of the Brunn--Minkowski and Pr\'ekopa--Leindler theorems,
including inequalities for log concave functions, and with an application
to the diffusion equation. \emph{Journal of Functional Analysis} 22
(1976), 366--389.
\bibitem{caffarelli} L. A. Caffarelli.
Monotonicity properties of optimal transportation and the FKG and related
inequalities. \emph{Communications in Mathematical Physics} 214 (2000),
547--563.
\bibitem{ledoux} M. Ledoux.
\emph{The Concentration of Measure Phenomenon}.
American Mathematical Society, 2001.
\end{thebibliography}
\endgroup

\appendix
\section{Compression of bounded-energy networks}\label{sec:norm-appendix}
The following deterministic result supplies the width-independent count
used in the strong-prior and $\gamma>1$ arguments. It depends on neither
Gaussianity nor any probabilistic property of the weights.
\begin{theorem}[Compression]\label{thm:compression}
Let $d\geq1$, $L\geq2$, and let $\phi:\mathbb R\to[-1,1]$ be
$1$-Lipschitz. Suppose a width-$N$ network $W$ satisfies
$\kappa_N(W)\leq K$, where $K\geq1$. Set
\begin{equation}\label{eq:compression-parameters}
 \begin{gathered}
 \eta=\frac{1}{12(L-1)},\qquad
 \rho=\eta 2^{-L}K^{-(L-1)/2},\\
 m_L(d,K):=M=\max\left\{\lceil9K\rceil,
 \left\lceil256\rho^{-4}\bigl(d+\log(L/\rho)\bigr)\right\rceil\right\}.
 \end{gathered}
\end{equation}
There is a width-$M$ network $\widetilde W$ such that
\begin{equation}\label{eq:compression-conclusion}
 \max_{x\in\{-1,1\}^d}
 |z_{\widetilde W}(x)-z_W(x)|\leq\frac1{12},
 \qquad \kappa_M(\widetilde W)\leq4K.
\end{equation}
If $W$ implements a Boolean function $g$ with strict margin $1/2$,
then multiplying the readout of $\widetilde W$ by $3/2$ gives a
width-$M$ implementation of $g$ with strict margin $1/2$,
weighted norm at most $9K$, and every weight of absolute value at most
$M$. In particular, $C_L(g)\leq M$. For fixed $L$, $M$ is bounded by a
polynomial in $d$ and $K$.
\end{theorem}

\begin{proof}
We first truncate large entries and then sample the hidden neurons.
All norms of rows below are Euclidean, and all matrix norms with
subscript $F$ are Frobenius norms.

\emph{Truncation.}
For $2\leq s\leq L$, let $r_i^s=\|W^s_{i,*}\|_2$. In each nonzero
row, delete entries larger in absolute value than
\[
 \tau_i^s=\frac{2r_i^s}{\rho\sqrt N}.
\]
Leave zero rows unchanged, write $\overline W^s$ for the resulting
matrices, and set $\overline W^1=W^1$. Here and below $W^L$ is the
single readout row. Because all hidden activations have absolute value
at most one, deleting a row's large entries changes its preactivation
by at most
\[
 \sum_{j:\,|W^s_{ij}|>\tau_i^s}|W^s_{ij}|
 \leq \frac{(r_i^s)^2}{\tau_i^s}
 =\frac\rho2\sqrt N\,r_i^s.
\]
The vector of changes at an internal layer therefore has norm at most
$(\rho/2)\sqrt N\|W^s\|_F$. Every original or truncated internal
matrix has operator norm at most $\sqrt K$, and the readout has norm
at most $\sqrt{K/N}$. The activation is a contraction in Euclidean
norm. Thus truncating layer $s<L$ changes the output by at most
\[
 \frac\rho2\sqrt{NK}\,K^{(L-s-1)/2}\sqrt{K/N}
 =\frac\rho2 K^{(L-s+1)/2}\leq\frac\eta2.
\]
Truncating the readout changes it by at most
$(\rho/2)\sqrt K\leq\eta/2$. These bounds also hold when other
layers have already been truncated. Telescoping over the $L-1$
truncations gives total output error at most $1/24$.

\emph{Independent samples and three events.}
For each hidden layer $s$, independently choose a list
$I_s(1),\ldots,I_s(M)$ of independent uniform indices in
$\{1,\ldots,N\}$. Repeated indices represent separate copies of a
neuron. Define the fully sampled network by
\begin{align*}
 \widetilde W^1_{a,*}&=W^1_{I_1(a),*},\\
 \widetilde W^s_{a,b}&=\frac NM
      \overline W^s_{I_s(a),I_{s-1}(b)}
      &&(2\leq s\leq L-1),\\
 \widetilde W^L_b&=\frac NM\overline W^L_{I_{L-1}(b)}.
\end{align*}
For the analysis expose these independent lists in reverse order,
$I_{L-1},\ldots,I_1$. At the stage sampling layer $m$, all earlier
layers still have their original, truncated weights. Let
$\overline h_j^m(x)$ denote their exactly computed activations. The
estimated sum into a retained row $i$ of layer $m+1$ is
\[
 \frac NM\sum_{b=1}^M
 \overline W^{m+1}_{i,I_m(b)}\overline h^m_{I_m(b)}(x).
\]
Its expectation is the corresponding full sum over $N$ neurons.
Conditional on the later lists, its summands are independent and
bounded in absolute value by $(N/M)\tau_i^{m+1}$. Hoeffding's inequality~\cite{hoeffding}
shows that its deviation exceeds
$(\rho/2)\sqrt N r_i^{m+1}$ with probability at most
\[
 2\exp\left(-\frac{\rho^4 M}{32}\right).
\]
A zero row is estimated exactly. The bound is uniform in the retained
row, so it remains valid after averaging over the later lists.

Let $A$ be the event that all these row estimates, at all stages and
all $2^d$ inputs, satisfy their deviation bounds. At the last hidden
layer there is one output row; at every other stage there are $M$
retained rows. Consequently
\begin{equation}\label{eq:compression-hoeffding}
 \Pr(A^c)\leq 2^{d+1}LM\exp(-\rho^4M/32)\leq\frac1{100}.
\end{equation}
For completeness, put $a=d+\log L+\log(1/\rho)$ and
$M_0=\lceil256\rho^{-4}a\rceil$. Then
$M_0\leq257\rho^{-4}a$, and $u\exp(-\rho^4u/32)$ decreases for
$u\geq M_0$. Substituting $M_0$ and using $\log a\leq a-1$ bounds
the logarithm of the right side before its final inequality by
\[
 -(7-\log2)d-6\log L-3\log(1/\rho)+\log514-1
 \leq-\log100.
\]
This also proves the bound when $M$ was enlarged to be at least $9K$.

Next define the total retained row mass using the \emph{original}
row norms:
\[
 R=N\sum_{s=2}^{L-1}\sum_{a=1}^M(r^s_{I_s(a)})^2
      +N\|W^L\|_2^2.
\]
Its expectation is
\[
 \mathbb E R
 =M\sum_{s=2}^{L-1}\|W^s\|_F^2+N\|W^L\|_2^2
 \leq MK.
\]
By Markov's inequality the event $B=\{R\leq8MK\}$ has probability
at least $7/8$. In particular, on $B$ the retained row mass at each
individual stage is at most $8MK$.

Finally, direct calculation gives
\begin{align*}
 \mathbb E\left[\frac dM\|\widetilde W^1\|_F^2\right]
     &=\frac dN\|W^1\|_F^2,\\
 \mathbb E\|\widetilde W^s\|_F^2
     &=\|\overline W^s\|_F^2 &&(2\leq s\leq L-1),\\
 \mathbb E\left[M\|\widetilde W^L\|_2^2\right]
     &=N\|\overline W^L\|_2^2.
\end{align*}
For an internal matrix, independence of the two adjacent layers'
index lists makes each index pair uniform on $\{1,\ldots,N\}^2$;
the factor $(N/M)^2$ is canceled by the $M^2/N^2$ sampling factor.
Independence among different matrix entries is unnecessary.
Thus $\mathbb E\kappa_M(\widetilde W)\leq K$, and the event
$C=\{\kappa_M(\widetilde W)\leq4K\}$ has probability at least $3/4$.
On $C$ every sampled internal matrix has operator norm at most
$2\sqrt K$, and the sampled readout has norm at most $2\sqrt{K/M}$.

The events need not be independent:
\[
 \Pr(A\cap B\cap C)\geq1-\frac1{100}-\frac18-\frac14>0.
\]
Fix a collection of lists in their intersection.

\emph{Reverse-order telescoping.}
Let $F_m$ be the network with hidden layers $m,\ldots,L-1$ sampled
and layers $1,\ldots,m-1$ still original and truncated. Incoming
weights to a sampled layer select its rows, and outgoing weights from
a sampled layer are multiplied by $N/M$. This defines the intermediate
networks with different widths in different layers. In particular,
$F_L$ is the truncated network and $F_1$ is the fully sampled network.

When comparing $F_{m+1}$ and $F_m$, layer $m$ computes the original
truncated neuron functions before selection. The only new discrepancy
is in the sums into the retained neurons of layer $m+1$.
For $m=L-1$, event $A$ bounds the output discrepancy directly by
$(\rho/2)\sqrt K\leq\eta/2$. For $m<L-1$, put $t=L-m$.
Events $A$ and $B$ bound the Euclidean discrepancy in layer $m+1$'s
preactivation by $(\rho/2)\sqrt{8MK}$. All matrices \emph{after}
layer $m+1$ already equal their final sampled matrices, so their
operator norms are controlled by $C$. Propagating the discrepancy
therefore changes the output by at most
\begin{align*}
 \frac\rho2\sqrt{8MK}\,(2\sqrt K)^{t-2}
       2\sqrt{K/M}
 &=\rho\,2^{t-1/2}K^{t/2}\\
 &\leq\eta\,2^{-3/2}\leq\frac\eta2,
 \qquad 2\leq t\leq L-1.
\end{align*}
There are $L-1$ sampling stages, contributing at most $1/24$ in total.
Adding the truncation error proves the approximation bound in
\eqref{eq:compression-conclusion}; event $C$ proves its norm bound.

If the original network has strict margin $1/2$, the sampled network
has strict margin $5/12$. Multiplying its linear readout by $3/2$
restores strict margin $1/2$ and multiplies the weighted norm by at
most $9/4$. At width $M$ the resulting entrywise bounds are
\[
 |W^1_{aj}|\leq\sqrt{9KM/d},\qquad
 |W^s_{ab}|\leq3\sqrt K\ (1<s<L),\qquad
 |W^L_b|\leq3\sqrt{K/M}.
\]
Each is at most $M$ because $d\geq1$ and $M\geq9K\geq9$.
The definition of $C_L$ now applies. Finally, substituting
\eqref{eq:compression-parameters} gives
$\rho^{-4}=12^4(L-1)^4 16^L K^{2(L-1)}$, which establishes the stated
polynomial dependence at fixed depth.
\end{proof}

% Self-contained proof fragment. Requires amsmath, amssymb, amsthm and
% theorem environments lemma/proposition. No references to older manuscripts.
\section{A finite family of entropy certificates}\label{sec:sparse-dual}
The following bound replaces a discretization of the full feature space by
sparse linear tests. Its logarithmic dependence on the number of observables
is essential when those observables are indexed by all pairs of Boolean
inputs.

Let $(\Omega,\mathcal A,\mu)$ be any probability space, let
$F=(F_1,\ldots,F_D):\Omega\to[-1,1]^D$ be measurable, and write
\[
 \Lambda(\lambda)=\log\int e^{\lambda\cdot F}\,d\mu,
 \qquad
 \Psi_G(q)=\max_{\|\lambda\|_1\leq G}
       \{\lambda\cdot q-\Lambda(\lambda)\},\qquad G>0.
\]
The maximum exists: its domain is compact and its objective is continuous.
The norm $\|\lambda\|_1$ here is the sum of the absolute values of its
$D$ coordinates.

\begin{lemma}[A bounded exponential tilt with controlled error]
\label{lem:bounded-dual-lift}
For every $q\in\mathbb R^D$, there is a maximizer $\lambda_*$ in the
preceding display such that the probability measure
\[
 \frac{d\nu_*}{d\mu}
     =\exp\{\lambda_*\cdot F-\Lambda(\lambda_*)\}
\]
satisfies the exact identity
\begin{equation}\label{eq:dual-lift-identity}
 \Psi_G(q)
   =D(\nu_*\Vert\mu)
      +G\left\|q-\int F\,d\nu_*\right\|_\infty.
\end{equation}
Consequently, if $\Psi_G(q)\leq K$, then
\[
 D(\nu_*\Vert\mu)\leq K,
 \qquad
 \left\|q-\int F\,d\nu_*\right\|_\infty\leq K/G.
\]
\end{lemma}
\begin{proof}
Boundedness of $F$ permits differentiation under the integral and gives
$\nabla\Lambda(\lambda_*)=\int F\,d\nu_*=:m_*$. The first-order
condition for a differentiable concave function maximized over the convex
$\ell^1$ ball is
\[
 (q-m_*)\cdot(\lambda-\lambda_*)\leq0
       \qquad(\|\lambda\|_1\leq G).
\]
Taking the maximum of the left linear functional over that ball gives
\[
 \lambda_*\cdot(q-m_*)=G\|q-m_*\|_\infty.
\]
Moreover,
$D(\nu_*\Vert\mu)=\lambda_*\cdot m_*-\Lambda(\lambda_*)$.
Adding these two identities proves \eqref{eq:dual-lift-identity}.
Relative entropy is nonnegative, which proves the stated consequences.
The log density of $\nu_*$ relative to $\mu$ is bounded, so its entropy
is finite. No support, density, or nonsingularity assumption on $\mu$
was used.
\end{proof}

For an integer $s\geq1$, introduce the finite set
\[
 \mathcal V_{G,s}
 =\left\{\frac1s\sum_{j=1}^s a_j:
        a_j\in\{0,Ge_1,-Ge_1,\ldots,Ge_D,-Ge_D\}\right\}.
\]
Thus $\#\mathcal V_{G,s}\leq(2D+1)^s$ and every vector in this set has
$\ell^1$ norm at most $G$.

\begin{lemma}[Sparse approximation of the dual value]
\label{lem:sparse-dual-net}
For every $q\in\mathbb R^D$,
\begin{equation}\label{eq:sparse-dual-net}
 \Psi_G(q)\leq\frac{G^2}{2s}
   +\max_{v\in\mathcal V_{G,s}}
        \{v\cdot q-\Lambda(v)\}.
\end{equation}
\end{lemma}
\begin{proof}
Fix $\lambda$ with $\|\lambda\|_1\leq G$. Let the random vector $A$
take the value $G\operatorname{sign}(\lambda_i)e_i$ with probability
$|\lambda_i|/G$ for each nonzero coordinate, and take value zero with the
remaining probability. Then $\mathbb EA=\lambda$.
Take independent copies $A_1,\ldots,A_s$ and put
$V=s^{-1}\sum_j A_j\in\mathcal V_{G,s}$.
For each fixed $\omega\in\Omega$, the scalar $A\cdot F(\omega)$ lies
in $[-G,G]$. Hoeffding's exponential-moment inequality therefore gives
\[
 \mathbb E\exp\{(V-\lambda)\cdot F(\omega)\}
 \leq\exp\{G^2/(2s)\}.
\]
Tonelli's theorem and Jensen's inequality imply
\begin{align*}
 \mathbb E\Lambda(V)
 &=\mathbb E\log\int e^{V\cdot F}\,d\mu\\
 &\leq\log\mathbb E\int e^{V\cdot F}\,d\mu\\
 &\leq\Lambda(\lambda)+G^2/(2s).
\end{align*}
Since $\mathbb E(V\cdot q)=\lambda\cdot q$, at least one realization
$v\in\mathcal V_{G,s}$ satisfies
\[
 v\cdot q-\Lambda(v)
 \geq\lambda\cdot q-\Lambda(\lambda)-G^2/(2s).
\]
Take $\lambda$ to be a maximizer defining $\Psi_G(q)$.
\end{proof}

\begin{proposition}[A nonasymptotic entropy certificate]
\label{prop:finite-entropy-certificate}
Let $X_1,\ldots,X_N$ be independent with law $\mu$, and put
$q_N=N^{-1}\sum_{i=1}^N F(X_i)$. For every $K>0$, $G>0$, and integer
$s\geq1$,
\begin{equation}\label{eq:finite-entropy-tail}
 \Pr\{\Psi_G(q_N)>K\}
 \leq\exp\left\{s\log(2D+1)
          -N\left(K-\frac{G^2}{2s}\right)\right\}.
\end{equation}
Outside this exceptional event, the empirical moment vector is within
$K/G$ in maximum norm of the moments of an exponential tilt with parameter
$\ell^1$ norm at most $G$ and relative entropy at most $K$.
In particular, choosing $s=\lceil G^2\rceil$ bounds the right side by
$\exp\{\lceil G^2\rceil\log(2D+1)-N(K-1/2)\}$.
\end{proposition}
\begin{proof}
For each fixed $v\in\mathcal V_{G,s}$, independence gives exactly
\[
 \mathbb E\exp\{N[v\cdot q_N-\Lambda(v)]\}=1.
\]
Markov's inequality bounds the probability that its bracket exceeds
$K-G^2/(2s)$ by $\exp\{-N(K-G^2/(2s))\}$.
The sparse approximation lemma and a union bound over
$\mathcal V_{G,s}$ prove \eqref{eq:finite-entropy-tail}.
The bounded-dual lift lemma gives the final assertion.
\end{proof}

\paragraph{Application to hidden-layer covariances.}
Take $m$ input points, $D=m^2$, $\mu=\mathcal N(0,Q)$ for any positive
semidefinite covariance $Q$, and $F_{ij}(z)=\tanh(z_i)\tanh(z_j)$,
indexed by all ordered pairs $(i,j)$. Reshape a maximizing parameter
$\lambda_*$ into a matrix and replace it by its symmetric part $\Gamma$.
This preserves its quadratic form and does not increase its entrywise
$\ell^1$ norm. The lifted covariance is then exactly
\[
 \mathsf F_\Gamma(Q)=
 \frac{\mathbb E_{Z\sim\mathcal N(0,Q)}
    [\tanh(Z)\tanh(Z)^\top e^{\tanh(Z)^\top\Gamma\tanh(Z)}]}
      {\mathbb E_{Z\sim\mathcal N(0,Q)}
    e^{\tanh(Z)^\top\Gamma\tanh(Z)}}.
\]
Thus, except with the probability in \eqref{eq:finite-entropy-tail}, the
empirical activation covariance $\widehat Q$ satisfies
\[
 \|\widehat Q-\mathsf F_\Gamma(Q)\|_{\max}\leq K/G,
 \qquad \|\Gamma\|_{1,\mathrm{ent}}\leq G.
\]
This assertion applies directly to Gaussian conditional row laws, uniformly
in their covariance, including singular covariances. For the actual
non-Gaussian rows in this paper, apply the distribution-free certificate
to their true law and then use the relative Lindeberg comparison
\eqref{eq:certificate-transfer}. On the Boolean cube,
$m=2^d$ and the logarithmic factor is
$\log(2m^2+1)\leq(2d+2)\log2$, which is linear in $d$.

\section{Covariance stability and finite realization}\label{sec:compression}
Write $\phi=\tanh$. Write $\|A\|_{\max}=\max_{ij}|A_{ij}|$ and
$\|A\|_{1,\mathrm{ent}}=\sum_{ij}|A_{ij}|$ for matrices on the finite input
set. These are different norms; their dual pairing is what avoids a factor
of $2^d$ below.
For a positive semidefinite covariance $Q$ and a symmetric matrix $\Gamma$,
let
\[
 \mathsf F_\Gamma(Q)=
 \frac{\mathbb E_{Z\sim\mathcal N(0,Q)}
 [\phi(Z)\phi(Z)^\top\exp\{\phi(Z)^\top\Gamma\phi(Z)\}]}
 {\mathbb E_{Z\sim\mathcal N(0,Q)}
 [\exp\{\phi(Z)^\top\Gamma\phi(Z)\}]}.
\]
All activations are coordinatewise. The definition is valid for singular $Q$.

\begin{lemma}[A dimension-independent covariance comparison]\label{lem:tilted-covariance}
If $g=\|\Gamma\|_{1,\mathrm{ent}}$, then for all positive semidefinite
$Q,Q'$ of the same size,
\[
 \|\mathsf F_\Gamma(Q)-\mathsf F_\Gamma(Q')\|_{\max}
 \leq 8(1+g)^2\|Q-Q'\|_{\max}.
\]
\end{lemma}
\begin{proof}
Put $V(z)=\phi(z)^\top\Gamma\phi(z)$ and
$Q_t=(1-t)Q+tQ'$. For a twice differentiable bounded test function $F$, let
$\mathbb E_t$ denote expectation under the Gaussian of covariance $Q_t$
reweighted by $e^V$. Gaussian differentiation and the quotient rule give
\begin{align*}
 \frac{d}{dt}\mathbb E_tF
 =\frac12\sum_{a,b}(Q'-Q)_{ab}\bigl\{
 &\mathbb E_t[\partial_{ab}F+
       \partial_aF\partial_bV+\partial_bF\partial_aV]\\
 &+\operatorname{Cov}_t(F,\partial_{ab}V+
                                   \partial_aV\partial_bV)\bigr\}.
\end{align*}
This formula has no reciprocal partition-function bound: all terms are
expectations under a probability measure. Define
$M_1(H)=\sup_z\sum_a|\partial_aH(z)|$ and
$M_2(H)=\sup_z\sum_{a,b}|\partial_{ab}H(z)|$.
Because $|\phi|,|\phi'|\leq1$ and $|\phi''|\leq2$,
\[
 M_1(V)\leq2g,\qquad M_2(V)\leq6g.
\]
For $F(z)=\phi(z_i)\phi(z_j)$, including $i=j$,
$\|F\|_\infty\leq1$, $M_1(F)\leq2$, and $M_2(F)\leq6$.
The absolute derivative is therefore at most
\[
 \frac12\|Q-Q'\|_{\max}
 \{6+8g+2(6g+4g^2)\}
 =(3+10g+4g^2)\|Q-Q'\|_{\max}.
\]
Integrating in $t$ proves the assertion, with the displayed slightly larger
constant. One may initially replace both covariances by $Q+\epsilon I$ and
$Q'+\epsilon I$. The functions and all derivatives used above are bounded,
and the partition functions lie between $e^{-g}$ and $e^g$, so weak Gaussian
convergence lets $\epsilon\downarrow0$. This also proves the singular case.
\end{proof}

\begin{proposition}[Finite realization of a bounded-tilt tower]\label{prop:tilt-compress}
Fix $L\geq2$, $K\geq1$, and $G\geq1$. Let $\mu_0$ be the actual
first-layer activation-field law, with independent incoming coordinates
$\xi_j/\sqrt d$ of law $\nu$. Suppose symmetric matrices $\Gamma_1,\ldots,\Gamma_{L-1}$ and a vector
$\beta\in\mathbb R^{\mathcal X}$ satisfy
\[
 \|\Gamma_s\|_{1,\mathrm{ent}}\leq G,\qquad
 \|\beta\|_1\leq4K.
\]
Define $Q_1$ by tilting $\mu_0$ with $e^{h^T\Gamma_1h}$,
and $Q_s=\mathsf F_{\Gamma_s}(Q_{s-1})$ for $s\geq2$.
Then the function $z(x)=(Q_{L-1}\beta)_x$ can be uniformly approximated to
accuracy $1/8$ by a depth-$L$ network whose width and weights are bounded by
$\operatorname{poly}_L(d,K,G)$.
If $f(x)z(x)\geq7/8$ everywhere, the finite network implements $f$ with margin at least $3/4$.
\end{proposition}
\begin{proof}
Let $b=8(1+G)^2$ and
\[
 \eta=[32K(L-1)b^{L-2}]^{-1}.
\]
Choose the common width $n$ to be the least integer satisfying
\[
 n\geq2\eta^{-2}\log\bigl(8(L-1)\,2^{2d}\bigr).
\]
The construction is random only to prove existence.
The matrices $\Gamma_s$ may be indexed by the entire cube; they select the
sampling laws, but are not stored or evaluated by the resulting finite
network. No efficient sampling or construction algorithm is claimed.
At the first layer, independently sample $n$ weight vectors from their
actual independent $\nu$ coordinate law divided by $\sqrt d$, tilted by
\[
 \exp\{h(w_i)^\top\Gamma_1h(w_i)\},\qquad
 h(w_i)_x=\phi(w_i\cdot\bar x).
\]
Their activation vectors have exactly the tilted law defining $Q_1$.
At each subsequent hidden layer, conditional on the already sampled previous
activation vectors $h_1,\ldots,h_n$, independently sample $n$ standard
Gaussian row vectors $\xi_j\in\mathbb R^n$ tilted by
\[
 \exp\{\phi(Z_j)^\top\Gamma_s\phi(Z_j)\},
 \qquad Z_j(x)=\frac1{\sqrt n}\sum_{i=1}^n\xi_{ji}h_i(x).
\]
Use the actual network weights $\xi_{ji}/\sqrt n$.
If $\widehat Q_{s-1}=n^{-1}\sum_i h_ih_i^\top$ is the preceding empirical
activation covariance, then, conditionally, $Z_j$ before tilting is Gaussian
with exactly that covariance. Thus the new activation vectors are
conditionally independent, with second moment
$\mathsf F_{\Gamma_s}(\widehat Q_{s-1})$.

Every covariance entry is the empirical mean of $n$ variables in $[-1,1]$.
Hoeffding's inequality, a union bound over $2^{2d}$ entries, and conditioning
at successive layers show, with probability at least $3/4$, that
\[
 \|\widehat Q_1-Q_1\|_{\max}\leq\eta,\qquad
 \|\widehat Q_s-\mathsf F_{\Gamma_s}(\widehat Q_{s-1})\|_{\max}\leq\eta
 \quad(2\leq s<L).
\]
Lemma~\ref{lem:tilted-covariance} then implies
\[
 \|\widehat Q_{L-1}-Q_{L-1}\|_{\max}
 \leq (L-1)b^{L-2}\eta.
\]
For each last-layer activation vector $h_i$, choose the deterministic
readout weight $v_i=(\beta\cdot h_i)/n$. The resulting output is
$\widehat z(x)=(\widehat Q_{L-1}\beta)_x$, and therefore
\[
 \|\widehat z-z\|_\infty
 \leq4K(L-1)b^{L-2}\eta=1/8.
\]
Although $\beta$ is indexed by the entire input cube, its only role in the
finite network is to select these $n$ real readout weights; each satisfies
$|v_i|\leq4K/n$. It is not an additional computational input to the network.

It remains to bound the sampled weights without losing polynomial size.
For every sampled row, its tilted density relative to its reference
is at most $e^{2G}$, since the quadratic tilt lies in $[-G,G]$.
Put $a_0=\min\{m/4,1/4\}$ and $B_0=\sqrt2$.
Equation~\eqref{eq:square-moment} and the Gaussian square moment show
that every raw reference coordinate obeys $\E e^{a_0\xi^2}\leq B_0$.
This also bounds the actual first-layer coordinates divided by $\sqrt d$.
Let $D_n=n(d+1)+(L-2)n^2$ and choose
\[
 B=\left[\frac{2G+\log(4B_0\max\{D_n,1\})}{a_0}\right]^{1/2}.
\]
The probability that any one coordinate exceeds $B$ in absolute value
is at most $B_0e^{2G-a_0B^2}$, also conditionally on previous layers.
A union bound gives probability at most $1/4$ that any incoming coordinate
exceeds $B$. The actual internal weights divide the raw coordinates by
$\sqrt n$ and so satisfy the same bound. Intersecting this event with the
covariance-approximation event leaves probability at least $1/2$.
Fix a realization in the intersection. Both $n$ and
$\max\{B,4K/n\}$ are bounded by
polynomials in $d,K,G$ at fixed depth. The stated margin follows from $7/8-1/8=3/4$.
\end{proof}

\begin{lemma}[Counting finite realizations]\label{lem:tilt-count}
Fix $L$, width $n$, and an entrywise bound $B\geq1$. The Boolean functions
having depth-$L$, width-$n$ implementations with margin at least $3/4$ and
weights bounded by $B$ number at most
\[
 \exp\{P\log(3+16B\Lambda)\},\qquad
 P=n(d+1)+(L-2)n^2+n,
 \quad \Lambda=(d+Ln)[1+n(B+1)]^{L-1}.
\]
\end{lemma}
\begin{proof}
Round every weight to a multiple of $(8\Lambda)^{-1}$.
The first hidden layer's coordinatewise error is at most $(d+1)/(8\Lambda)$.
Every subsequent hidden layer has error at most $n(B+1)$ times the preceding
error plus $n/(8\Lambda)$; the same estimate applies at the linear readout.
Induction bounds the final error by $1/8$. Each rounded coordinate has at
most $3+16B\Lambda$ possible values, so the $P$-parameter grid has the stated
size. Rounding preserves every represented Boolean sign function.
\end{proof}



\end{document}
